(1) 磁通链接分别为 \(\Phi_1=LI_1+MI_2\)、\(\Phi_2=MI_1+LI_2\)。Kirchhoff 方程为
\[
L\dot I_1+M\dot I_2+RI_1=0,\qquad
M\dot I_1+L\dot I_2+RI_2=0.
\]

(2) 定义 \(I_+=I_1+I_2,\ I_-=I_1-I_2\)，相加减得
\[
(L+M)\dot I_++RI_+=0,\quad
(L-M)\dot I_-+RI_-=0.
\]
因此 \(I_\pm(t)=I_\pm(0)e^{-t/\tau_\pm}\)，其中 \(\tau_\pm=(L\pm M)/R>0\)，回路电流由 \(I_{1,2}=(I_+\pm I_-)/2\) 恢复。

(3) 磁能 \(U=\frac12L(I_1^2+I_2^2)+MI_1I_2=\frac14[(L+M)I_+^2+(L-M)I_-^2]\)。因 |M|<L 两个系数严格正，故非零电流必有正能量。由电路方程乘以对应电流并相加得
\[
\dot U=-R(I_1^2+I_2^2)\le0.
\]
这就是电阻焦耳热守恒关系。

(4) M→0 时两个时间常数均趋 L/R，回路独立指数衰减。M→L⁻ 时反相模 \(\tau_-\to0\)，磁矩阵趋向奇异；严格 M=L 时失去反相磁能的正定性，因此不能不加限制地把 M=L 代入普通微分方程。
